\section{概览}
\begin{enumerate}
  \item 分数占比最高。
  \item 均、直接、只有、全部。
  \item 备考方法。
        \begin{enumerate}
          \item 体系化。
          \item 图示化，民商事\zd{法律关系一般都在 3 个以上当事人，得把两两当事人的法律关系搞清楚}。
          \item 案例化。
          \item 只要跟你正常的朴素的法律常识不认识一样，才去单独学习。
          \item 1 部民法典 + 3 个司法解释（总则编，担保制度解释，民间借贷合同）。
        \end{enumerate}
  \item 民法典。7 编。
        \begin{enumerate}
          \item 总则。10\%
                \begin{enumerate}
                  \item 民事法律关系。民事主体（自然人、法人、非法人）。
                  \item 民事法律事实。民事法律行为、代理、诉讼时效。
                \end{enumerate}
          \item 财产权。
                \begin{enumerate}
                  \item 物权。所有权；用益物权；担保物权；准物权/占有。
                  \item 合同。通则；19 类典型；准合同/无因管理 + 不当得利。
                \end{enumerate}
          \item 人身权。
                \begin{enumerate}
                  \item 人格权。2 类一般人格权；9 类具体人格权；个人信息保护。
                  \item 婚姻家庭。结婚；家庭关系；离婚；收养。
                \end{enumerate}
          \item 人身 + 财产（综合性权利）。
                \begin{enumerate}
                  \item 继承。法定继承 vs 遗嘱继承；代位 vs 转；遗赠 vs 遗赠抚养。
                \end{enumerate}
          \item 侵权责任。共同侵权；精神损害赔偿；7 类特殊主体责任；7 类典型侵权责任。
        \end{enumerate}
\end{enumerate}

\section{民事法律关系的基本原理}
\begin{enumerate}
  \item \zd{产生私法效果（人身、财产）\yjt 法律约束力}。
  \item 考点，哪些归民法管 \yjt 民法的调整对象。
        \zd{调整平等主体的自然人、法人、非法人组织之间的人身关系和财产关系}。
        \begin{enumerate}
          \item 行政、生活行为。\al{交警贴罚单；公安局买菜；恋爱、同居、朋友关系}。
          \item 好意施惠。
                \begin{enumerate}
                  \item \al{邀请看电影；火车到站叫醒；志愿者}。
                  \item \al{请客吃饭极力劝酒}，有过错的侵权行为，一般侵权。\al{酒后照顾义务}，赔偿 \hz 补偿。
                  \item \al{搭便车驾驶人与第三人共同过错}，按份责任。没考过。
                \end{enumerate}
          \item 戏谑行为。\al{找到错别字给你 10W}，得能产生合理信赖。
        \end{enumerate}
  \item 3 大要素。
        \begin{enumerate}
          \item 民事主体。自然人、法人、非法人组织；国家或集体，\al{国家发行国债；无主物遗失物}。
          \item 民事内容。\yi 民事权利，依法为或不为的自由；\er 民事义务，当为或不为某种行为的约束。
          \item 民事客体。权利义务所指向的对象。
                \begin{enumerate}
                  \item \zd{物权} \yjt 物为原则，\zd{权利}为例外；
                  \item \zd{债权} \yjt 行为（特定当事人之间请求为一定给付），作为 \hz \zd{不作为}。\al{竞业禁止协议}。
                  \item 人身权 \yjt 人身利益；人格权 \yjt 人格利益。
                  \item 身份权 \yjt 身份利益；
                  \item 知识产权 \yjt 智力成果；
                  \item 继承权 \yjt 遗产。
                \end{enumerate}
        \end{enumerate}
\end{enumerate}
